% !TeX root = ../教材题库_全部条目_转写工作稿.tex
% Source: XJ-B1-C05-TH-P158-Y1003
% 第5章 / 5.1 任意角和弧度制 / 教材 p.158、159 / PDF 第165、166页
% OCR未校对：公式与文字均可能有误；未完成数学TeX转写；禁止直接发给学生
% 原题及相邻原解截图（核对用）：../题目原图/XJ-B1-C05-TH-P158-Y1003-01.webp, ../题目原图/XJ-B1-C05-TH-P158-Y1003-02.webp
% 本文件只在转写完成后作为题干使用；OCR 条目尚不是可靠题干。
\begin{bankitem}{p.158 思考（含紧邻说明）}{第5章 / 5.1 任意角和弧度制 / 教材 p.158、159 / PDF 第165、166页}{XJ-B1-C05-TH-P158-Y1003}{OCR未校对：公式与文字均可能有误；未完成数学TeX转写；禁止直接发给学生}
\begin{ocrdraft}
[教材p.158 / PDF第165页]
\par
？思考如图5，1一3，在半径分别为0，0的任意两个不同
\par
的同心圆中，圆心角a所对应的弧长分别为厶，/2．分别计算
\par
由此你能发现什么？
\par
设a一刀。，根据初中所学知识可知
\par
图5．1一3
\par
刀丌厂1
\par
刀丌0
\par
180
\par
180\ensuremath{\cdot}
\par
因此
\par
77
\par
180．
\par
可以发现，当圆心角一定时，弧长与半径的比值一与半径无关，仅与圆心角仅
\par
的大小有关．由此启发我们可以利用圆的弧长与半径的比值去度量圆心角，从而找
\par
到一种新的度量角的方法．
\par
2丌厂
\par
例如，当仅一360。时，
\par
=2丌．同理可得，当a一180。时，
\par
厂
\par
厂
\par
90。时，
\par
当仅
\par
一60。时，
\par
[教材p.159 / PDF第166页]
\par
由于当I一厂时，
\par
=1，于是数学上把长度等于半径长的弧
\par
1rad
\par
所对的圆心角叫作1弧度的角，“弧度”用符号rad表示．如图
\par
0
\par
，4
\par
5．1一4，圆0的半径为。，的长等于。则／，AOB就是1弧度
\par
的角．
\par
图5．1一4
\par
如上所述，这样规定的1弧度角的大小是完全确定的，与所
\par
在圆的大小（即半径）无关．这种以“弧度”为单位来度量角的单位制叫作弧度制．
\par
在弧度制下，我们实现了用单位长度来度量角的大小．
\par
一般地，若半径为厂的圆心角所对的弧长为/，则弧长与半径的比值就是角
\par
的弧度数的绝对值，即
\par
这里，a的正负由角a的终边的旋转方向决定，即逆时针旋转为正，顺时针旋转为
\par
负．显然，正角的弧度数是一个正数，负角的弧度数是一个负数，零角的弧度数
\par
特别地，当角所在圆的半径是1（即单位圆）时，那么角a的弧度数的绝对值
\par
=/．也就是说，我们可以用单位圆中的弧长/来度量它所对的角a的弧度数的绝
\par
对值．
\end{ocrdraft}
\end{bankitem}
